\documentclass{situationist-preprint}

\title{Relativity of Simultaneity for Philosophers}
\msid{TSM-16-b}
\msfamily{TSM-16}
\msdiscipline{Time, Space and Modality}
\mstype{Companion manuscript}
\mskeywords{relativity of simultaneity; special relativity; conventionality; Malament; objective now}

\begin{document}
\maketitle

\begin{abstract}
TSM-16-a proves that Minkowski spacetime admits no simultaneity relation invariant under its symmetries. This note sets out the physics the proof relies on in elementary form: how simultaneity is defined by light signals, why observers in relative motion disagree, why the disagreement is not an artefact of convention, and what Malament's theorem adds. The upshot is that each inertial observer has a well-defined simultaneity relation, that no one of them is singled out by the spacetime, and that the present is a local notion. Simultaneity is frame-relative, and there is no objective now.
\end{abstract}

\section{Simultaneity by light signals}

\citet{Einstein1905} defined simultaneity for distant events operationally. An observer at rest sends a light signal at time $t_1$ by her clock to a distant mirror, and receives the reflection at $t_2$. The reflection event is assigned the time $\tfrac{1}{2}(t_1 + t_2)$, on the assumption that light takes as long to go as to return. Two events are simultaneous for her if they receive the same time.

Now let a second observer move with speed $v$ relative to the first, and apply the same procedure. Since the speed of light is the same for both, the midpoint rule assigns, in the second observer's frame, equal times to pairs of events that the first assigns different times. Quantitatively, events at distance $x$ along the line of motion that the first observer counts as simultaneous differ, for the second, by $\gamma v x / c^2$, where $\gamma = (1 - v^2/c^2)^{-1/2}$. This is the content of the boost lemma (Lemma~3.1 of TSM-16-a): for any two spacelike separated events there is an inertial frame in which they are simultaneous, and frames in which either precedes the other.

\section{Convention and fact}

\citet{Reichenbach1958} held that the choice of the midpoint is a convention: the one-way speed of light cannot be measured without a prior synchronization, so any value of the factor $\tfrac{1}{2}$ between $0$ and $1$ is admissible. If so, frame-relativity might seem to be an artefact of a conventional choice rather than a fact about spacetime.

\citet{Malament1977} answered that, given an inertial observer's worldline, the standard relation is the only nontrivial equivalence relation definable from causal connectibility and that worldline. Standard simultaneity is therefore not conventional relative to an observer. Malament's theorem and the result of TSM-16-a are complementary. The first shows that each inertial observer has a privileged simultaneity relation; the second, Theorem~3.2 of TSM-16-a, shows that no simultaneity relation is privileged without reference to an observer, since none is invariant under the boosts that carry one inertial worldline to another. Simultaneity is a relation between two events and a frame.

\begin{objection}
A theory with a privileged rest frame, undetectable but real, is empirically equivalent to special relativity and restores absolute simultaneity.
\end{objection}

\begin{reply}
TSM-16-a answers this neo-Lorentzian proposal in its Objection~1: the privileged frame is an empirically equivalent rival without independent motivation, since the only motive offered for it is the A-theorist's demand for a global present. By the account of underdetermination of SCI-04, such rivals do not threaten belief in the theory they mimic.
\end{reply}

\begin{objection}
If simultaneity is frame-relative, then by the argument of \citet{Putnam1967} future events are already real for some observer, and the future is fixed.
\end{objection}

\begin{reply}
The argument chains simultaneities belonging to different frames, and such chains relate events that are timelike separated; TSM-16-a answers it in its Objection~3 by defining what is settled for an event through its causal past rather than through simultaneity. On that account frame-relative simultaneity has no consequence for the openness of the future, which concerns settledness relative to a moment (TSM-12).
\end{reply}

\section{What remains of the present}

Each observer has a local present: the events that every frame dates within a short interval of her proper time, which by Proposition~4.2 of TSM-16-a fill a region whose extent is bounded by the length of the interval. As the interval shrinks, so does the region, and at the limit the present is the here-now. That is all that relativity leaves of the present, and it is enough for the experience and practice that A-theorists cite, which are local.

\begin{thebibliography}{99}\small
\bibitem[Einstein(1905)]{Einstein1905} Einstein, A. (1905). Zur Elektrodynamik bewegter K\"orper. \emph{Annalen der Physik}, 17, 891--921.
\bibitem[Malament(1977)]{Malament1977} Malament, D. (1977). Causal theories of time and the conventionality of simultaneity. \emph{No\^us}, 11(3), 293--300.
\bibitem[Putnam(1967)]{Putnam1967} Putnam, H. (1967). Time and physical geometry. \emph{The Journal of Philosophy}, 64(8), 240--247.
\bibitem[Reichenbach(1958)]{Reichenbach1958} Reichenbach, H. (1958). \emph{The Philosophy of Space and Time}. Dover. (Original work published 1928.)
\end{thebibliography}

\end{document}
